\documentclass[11pt]{article}

\usepackage[margin=1in]{geometry}
\usepackage{amsmath}
\usepackage{amssymb}
\usepackage{booktabs}
\usepackage[hidelinks]{hyperref}

\newcommand{\softmax}{\operatorname{softmax}}

\title{\textbf{Exercise 06: Multi-Query and Grouped-Query Attention}}
\author{}
\date{}

\begin{document}
\maketitle

\begin{center}
\emph{How multiple query heads can share fewer key/value heads.}
\end{center}

\section*{1. MHA, GQA, and MQA}

Let $H_q$ be the number of query heads and $H_{kv}$ the number of key/value heads.
\[
\boxed{\text{MHA: }H_{kv}=H_q}
\]
\[
\boxed{\text{GQA: }1<H_{kv}<H_q}
\]
\[
\boxed{\text{MQA: }H_{kv}=1}
\]
The number of query heads does not decrease when $H_{kv}$ decreases.

For $H_q=4$ and GQA with $H_{kv}=2$, the group size is
\[
G=\frac{H_q}{H_{kv}}=2.
\]
Using zero-based indices,
\[
g(h)=\left\lfloor\frac{h}{G}\right\rfloor,
\]
so the mapping is
\[
\boxed{[0,0,1,1]}.
\]
For MQA, every query head maps to the single K/V head:
\[
\boxed{[0,0,0,0]}.
\]

\section*{2. Shared K/V attention}

For query head $h$,
\[
\boxed{
\mathrm{head}_h
=\softmax\left(
\frac{Q_hK_{g(h)}^\top}{\sqrt{d_k}}
\right)V_{g(h)}
}.
\]
$Q_h$ remains head-specific even when K/V is shared.

Use
\[
q_1=\begin{bmatrix}1&0\end{bmatrix},
\qquad
q_2=\begin{bmatrix}0&1\end{bmatrix},
\]
\[
K_{\mathrm{shared}}=
\begin{bmatrix}
1&0\\
0&1
\end{bmatrix},
\qquad
V_{\mathrm{shared}}=
\begin{bmatrix}
1&0\\
1&1
\end{bmatrix},
\qquad d_k=2.
\]
The raw scores are
\[
q_1K_{\mathrm{shared}}^\top=\begin{bmatrix}1&0\end{bmatrix},
\qquad
q_2K_{\mathrm{shared}}^\top=\begin{bmatrix}0&1\end{bmatrix}.
\]
After scaling by $\sqrt{2}$ and applying softmax,
\[
A_1\approx\begin{bmatrix}0.6698&0.3302\end{bmatrix},
\qquad
A_2\approx\begin{bmatrix}0.3302&0.6698\end{bmatrix}.
\]
Therefore
\[
O_1=A_1V_{\mathrm{shared}}
\approx\begin{bmatrix}1.0000&0.3302\end{bmatrix},
\]
\[
O_2=A_2V_{\mathrm{shared}}
\approx\begin{bmatrix}1.0000&0.6698\end{bmatrix}.
\]
The same K/V state can therefore yield different attention patterns and outputs because the queries are different.

\section*{3. Projection widths}

Using row-vector mathematical notation,
\[
W_Q\in\mathbb{R}^{d_{\mathrm{model}}\times(H_qd_k)},
\]
\[
W_K\in\mathbb{R}^{d_{\mathrm{model}}\times(H_{kv}d_k)},
\qquad
W_V\in\mathbb{R}^{d_{\mathrm{model}}\times(H_{kv}d_v)},
\]
\[
W_O\in\mathbb{R}^{(H_qd_v)\times d_{\mathrm{model}}}.
\]
For $d_{\mathrm{model}}=8$, $H_q=4$, and $d_k=d_v=2$:

\begin{center}
\begin{tabular}{@{}lccc@{}}
\toprule
Architecture & Q width & K width & V width\\
\midrule
MHA ($H_{kv}=4$) & 8 & 8 & 8\\
GQA ($H_{kv}=2$) & 8 & 4 & 4\\
MQA ($H_{kv}=1$) & 8 & 2 & 2\\
\bottomrule
\end{tabular}
\end{center}

The combined bias-free K/V projection parameter count is
\[
\boxed{d_{\mathrm{model}}H_{kv}(d_k+d_v)}.
\]
Thus
\[
\text{MHA}=128,\qquad
\text{GQA}=64,\qquad
\text{MQA}=32.
\]
$W_Q$ and $W_O$ remain $8\times8$ in this toy setup.

\section*{4. GQA tensor layout}

A practical layout is
\[
Q:[B,H_q,L,D],
\qquad
K,V:[B,H_{kv},L,D].
\]
For $B=1$, $H_q=4$, $H_{kv}=2$, $L=3$, and $D=2$:
\[
Q:[1,4,3,2],
\qquad
K,V:[1,2,3,2].
\]
The attention score and output shapes are
\[
\text{scores}:[1,4,3,3],
\qquad
\text{output}:[1,4,3,2].
\]

A teaching implementation may materialize query-head-shaped K/V tensors such as
\[
K_{\mathrm{for\ q}},V_{\mathrm{for\ q}}:[B,H_q,L,D]
\]
by indexing each query head's assigned K/V head. This does \emph{not} imply that the physical KV cache must store $H_q$ copies. Optimized GQA implementations can map query heads to shared K/V heads without duplicating the stored cache.

\section*{5. KV-cache size}

For one layer,
\[
K_{\mathrm{cache}}\in\mathbb{R}^{B\times H_{kv}\times L\times d_k},
\qquad
V_{\mathrm{cache}}\in\mathbb{R}^{B\times H_{kv}\times L\times d_v}.
\]
The cached scalar-element count is
\[
\boxed{BLH_{kv}(d_k+d_v)}.
\]
Across $N_{\mathrm{layers}}$ layers,
\[
\boxed{BN_{\mathrm{layers}}LH_{kv}(d_k+d_v)}.
\]
For $B=1$, one layer, $L=3$, and $d_k=d_v=2$:
\[
\text{MHA}=48,\qquad
\text{GQA}=24,\qquad
\text{MQA}=12.
\]
The cache ratio is therefore $4:2:1$ for this setup.

\section*{6. What actually gets cheaper?}

Reducing $H_{kv}$ directly reduces:
\begin{itemize}
    \item K projection width and parameters,
    \item V projection width and parameters,
    \item KV-cache storage,
    \item K/V memory traffic during decoding.
\end{itemize}

However, $H_q$ remains fixed. The model still has the same number of query heads, and each query head still forms its own attention distribution over sequence positions. Therefore
\[
\boxed{
H_{kv}\downarrow
\not\Rightarrow
\text{all attention computation shrinks by the same factor}
}.
\]

\section*{Final answer}

MQA can use one shared K/V head while retaining multiple distinct query-head attention patterns because the query projection still produces $H_q$ distinct query heads. Those different $Q_h$ vectors interact with the same K/V state to produce different score vectors, softmax distributions, and outputs.

\section*{Bonus}

KV-cache storage is linear in $L$, so doubling sequence length doubles cache size for MHA, GQA, and MQA. Their relative ratios do not change; those ratios are determined by $H_{kv}$.

\section*{Numerical verification}

The companion \texttt{answer\_key.py} verifies the head mappings, shared-K/V numerical example, projection widths, cache ratios, and a GQA tensor attention calculation. It also compares a materialized vectorized teaching implementation with the canonical per-query-head formula.

\end{document}
